\documentclass[11pt,a4paper]{article}
\usepackage[utf8]{inputenc}
\usepackage[T1]{fontenc}
\usepackage[english]{babel}
\usepackage{amsmath,amssymb,amsthm}
\usepackage{geometry}
\usepackage{booktabs}
\usepackage{longtable}
\usepackage{xcolor}
\usepackage{hyperref}
\usepackage{enumitem}
\usepackage{fancyhdr}
\usepackage{titlesec}
\usepackage{listings}
\usepackage{algorithm}
\usepackage{algpseudocode}
\usepackage{tikz}
\usepackage{pgfplots}
\usepackage{graphicx}
\usepackage{caption}
\usepackage{subcaption}
\usepackage{float}
\usepackage{multirow}
\usepackage{array}
\usepackage{makecell}
\usepackage{fontawesome5}

\geometry{margin=2.5cm}
\pagestyle{fancy}
\fancyhf{}
\fancyhead[L]{\leftmark}
\fancyhead[R]{\thepage}
\renewcommand{\headrulewidth}{0.4pt}

\titleformat{\section}{\Large\bfseries}{\thesection}{1em}{}
\titleformat{\subsection}{\large\bfseries}{\thesubsection}{1em}{}
\titleformat{\subsubsection}{\normalsize\bfseries}{\thesubsubsection}{1em}{}

\definecolor{codegreen}{rgb}{0,0.6,0}
\definecolor{codegray}{rgb}{0.5,0.5,0.5}
\definecolor{codepurple}{rgb}{0.58,0,0.82}
\definecolor{backcolour}{rgb}{0.95,0.95,0.92}
\definecolor{darkblue}{rgb}{0,0,0.54}

\lstdefinestyle{mystyle}{
    backgroundcolor=\color{backcolour},
    commentstyle=\color{codegreen},
    keywordstyle=\color{blue},
    numberstyle=\tiny\color{codegray},
    stringstyle=\color{codepurple},
    basicstyle=\ttfamily\footnotesize,
    breakatwhitespace=false,
    breaklines=true,
    captionpos=b,
    keepspaces=true,
    numbers=left,
    numbersep=5pt,
    showspaces=false,
    showstringspaces=false,
    showtabs=false,
    tabsize=2
}
\lstset{style=mystyle}

\hypersetup{
    colorlinks=true,
    linkcolor=darkblue,
    filecolor=magenta,
    urlcolor=cyan,
    citecolor=green,
}

\title{\textbf{QPU-Estimator}\\\Large A Python Framework for IBM Quantum Hardware Resource Estimation\\\vspace{0.5em}\large Technical Documentation v0.1.0}
\author{Daniel Mart\'{i}n P\'{e}rez\\Quantum Computing Research Group\\Universidad Pablo de Olavides}
\date{\today}

\begin{document}
\maketitle
\tableofcontents
\newpage

%=============================================================================
\section{Introduction}
%=============================================================================

QPU-Estimator is a Python framework for estimating IBM Quantum hardware resource usage from arbitrary quantum circuits. Given a \texttt{QuantumCircuit} and a target IBM backend, it predicts execution time, optimal shot counts, accumulated error, and credit cost---before any real execution.

\subsection{Motivation}

IBM Quantum Runtime provides execution APIs but no upfront resource estimator. For quantum machine learning (QML) research---where backend choice justification, transpilation overhead comparison, and experiment budgeting are critical---this framework fills the gap.

\subsection{Key Features}

\begin{itemize}[leftmargin=2em]
    \item \textbf{Live backend profiling} via \texttt{qiskit-ibm-runtime}
    \item \textbf{Real transpilation} with \texttt{qiskit.transpiler}
    \item \textbf{Noise-aware fidelity} prediction (depolarizing + thermal relaxation + readout)
    \item \textbf{Shot optimization} using Hoeffding / Chernoff / Clopper-Pearson bounds
    \item \textbf{QML specialization} for VQC training, MAML/QMAML, and transfer learning
    \item \textbf{Training time estimation} including epochs, batches, queue time, and session mode
    \item \textbf{Export utilities} for LaTeX, CSV, and Markdown
    \item \textbf{Plugin system} for custom backends (IonQ example included)
\end{itemize}

\subsection{Architecture Overview}

\begin{lstlisting}[language=TeX,caption=Package Structure]
qpu_estimator/
|-- analyzer.py          # Circuit analysis (gate counts, depth, measurements)
|-- profiler.py          # Backend profiling (live IBM + mock fallback)
|-- transpiler.py        # Transpilation estimation (real + heuristic)
|-- estimator.py         # Resource estimation (time, shots, fidelity, credits)
|-- noise_model.py       # Noise-aware fidelity model
|-- shot_optimizer.py    # Statistical shot optimization
|-- qml_estimator.py     # QML-specific estimators
|-- training_time_estimator.py  # Full training time estimation
|-- export.py            # LaTeX/CSV/Markdown export
|-- plugin.py            # Plugin system for custom backends
|-- dashboard.py         # Simple web dashboard
|-- cli.py               # Command-line interface
|-- models.py            # Data models
\end{lstlisting}

%=============================================================================
\section{Installation}
%=============================================================================

\begin{lstlisting}[language=bash,caption=Installation from Source]
git clone https://github.com/danimap27/qpu-estimator.git
cd qpu-estimator
python -m venv .venv
source .venv/bin/activate
pip install -e ".[dev]"
\end{lstlisting}

\subsection{Dependencies}

\begin{table}[htbp]
\centering
\caption{Core Dependencies}
\begin{tabular}{@{}lll@{}}
\toprule
\textbf{Package} & \textbf{Version} & \textbf{Purpose} \\
\midrule
qiskit & $\geq$1.0 & Quantum circuit construction \\
qiskit-ibm-runtime & $\geq$0.25 & IBM Quantum backend access \\
numpy & $\geq$1.24 & Numerical computation \\
pytest & $\geq$7.0 & Testing (dev) \\
ruff & latest & Linting (dev) \\
mypy & latest & Type checking (dev) \\
\bottomrule
\end{tabular}
\end{table}

%=============================================================================
\section{Core API}
%=============================================================================

\subsection{QPUEstimator}

The main orchestrator class.

\begin{lstlisting}[language=Python,caption=Basic Usage]
from qiskit import QuantumCircuit
from qpu_estimator import QPUEstimator

circuit = QuantumCircuit(2)
circuit.h(0)
circuit.cx(0, 1)
circuit.measure_all()

estimator = QPUEstimator(
    use_live=True,           # Fetch live IBM data
    use_real_transpiler=True, # Use qiskit.transpiler
    ibm_token="YOUR_TOKEN"   # Optional: explicit token
)

report = estimator.estimate(circuit, backend_name="ibm_kingston")
print(f"Fidelity: {report.estimated_fidelity:.4f}")
print(f"Shots: {report.optimal_shots}")
print(f"Time: {report.estimated_execution_time_ms:.4f} ms")
\end{lstlisting}

\subsection{Configuration Options}

\begin{table}[htbp]
\centering
\caption{QPUEstimator Constructor Parameters}
\begin{tabular}{@{}lllp{5cm}@{}}
\toprule
\textbf{Parameter} & \textbf{Type} & \textbf{Default} & \textbf{Description} \\
\midrule
\texttt{config} & EstimationConfig & None & Custom estimation parameters \\
\texttt{use\_live} & bool & True & Fetch live IBM backend data \\
\texttt{ibm\_token} & str & None & IBM Quantum API token \\
\texttt{use\_real\_transpiler} & bool & True & Use qiskit.transpiler vs heuristic \\
\bottomrule
\end{tabular}
\end{table}

\subsection{EstimationReport}

The output data model.

\begin{lstlisting}[language=Python,caption=EstimationReport Fields]
@dataclass(frozen=True)
class EstimationReport:
    backend_name: str
    circuit_profile: CircuitProfile      # Structural analysis
    transpiled_depth: int                # After transpilation
    estimated_execution_time_ms: float
    optimal_shots: int
    estimated_fidelity: float
    estimated_credits: float
    swap_count: int
    notes: list[str]                     # Warnings and info
\end{lstlisting}

\subsection{Multi-Backend Comparison}

\begin{lstlisting}[language=Python,caption=Comparing Multiple Backends]
reports = estimator.compare_backends(
    circuit,
    ["ibm_fez", "ibm_marrakesh", "ibm_kingston"]
)
# Sorted by estimated fidelity (descending)
for r in reports:
    print(f"{r.backend_name}: fidelity={r.estimated_fidelity:.4f}")
\end{lstlisting}

%=============================================================================
\section{Backend Profiling}
%=============================================================================

\subsection{Live Profiling}

Fetches real calibration data from IBM Quantum:

\begin{itemize}[leftmargin=2em]
    \item T1/T2 coherence times per qubit
    \item Single-qubit gate errors
    \item Two-qubit gate errors per edge
    \item Readout errors
    \item Coupling map topology
    \item Basis gates
\end{itemize}

\begin{lstlisting}[language=Python,caption=Live Backend Profiling]
from qpu_estimator import QPUEstimator

estimator = QPUEstimator(use_live=True)
backends = estimator.list_live_backends()
print(f"Available: {backends}")
# ['ibm_fez', 'ibm_marrakesh', 'ibm_kingston']
\end{lstlisting}

\subsection{Mock Profiles (Offline)}

When IBM is unavailable, falls back to mock profiles:

\begin{table}[htbp]
\centering
\caption{Mock Backend Profiles}
\begin{tabular}{@{}llll@{}}
\toprule
\textbf{Backend} & \textbf{Qubits} & \textbf{Basis Gates} & \textbf{Topology} \\
\midrule
ibm\_heron & 133 & rz, sx, x, ecr & Heavy-hex \\
ibm\_brisbane & 127 & rz, sx, x, ecr & Heavy-hex \\
ibm\_sherbrooke & 127 & rz, sx, x, ecr & Heavy-hex \\
\bottomrule
\end{tabular}
\end{table}

%=============================================================================
\section{Transpilation Estimation}
%=============================================================================

\subsection{Real Transpilation}

Uses \texttt{generate\_preset\_pass\_manager} from Qiskit:

\begin{lstlisting}[language=Python,caption=Real Transpilation]
from qiskit.transpiler.preset_passmanagers import generate_preset_pass_manager

pm = generate_preset_pass_manager(
    optimization_level=2,
    coupling_map=backend_profile.coupling_map,
    basis_gates=backend_profile.basis_gates,
)
transpiled = pm.run(circuit)
\end{lstlisting}

\subsection{Heuristic Fallback}

When real transpilation fails, uses a swap insertion heuristic:

\begin{equation}
\text{swaps} = \lceil \text{CNOTs} \times \text{avg\_distance} \times 2 \rceil
\end{equation}

where \texttt{avg\_distance} approximates the average shortest path on the coupling map.

%=============================================================================
\section{Noise-Aware Fidelity}
%=============================================================================

\subsection{Noise Channels}

The noise model combines three channels:

\begin{enumerate}[leftmargin=2em]
    \item \textbf{Depolarizing:} From gate error rates
    \item \textbf{Thermal relaxation:} From T1/T2 decay during gate times
    \item \textbf{Readout:} From measurement error rates
\end{enumerate}

\subsection{Thermal Relaxation Model}

Per-layer fidelity using empirically calibrated exponential decay:

\begin{equation}
F_{\text{thermal}} = \prod_{q=1}^{n} \left( e^{-\frac{t_g}{2 T_1^{(q)}}} \cdot e^{-\frac{t_g}{2 T_{\phi}^{(q)}}} \right)^{d}
\end{equation}

where $t_g$ is average gate time, $T_1$ is relaxation time, $T_{\phi}$ is pure dephasing time, and $d$ is circuit depth.

\subsection{Configuration}

\begin{lstlisting}[language=Python,caption=Noise Model Configuration]
from qpu_estimator.estimator import EstimationConfig

config = EstimationConfig(
    use_noise_aware_fidelity=True,
    target_precision=0.01,
    confidence=0.95,
)
\end{lstlisting}

%=============================================================================
\section{Shot Optimization}
%=============================================================================

\subsection{Statistical Bounds}

Three methods for shot count optimization:

\begin{table}[htbp]
\centering
\caption{Shot Optimization Methods}
\begin{tabular}{@{}llp{6cm}@{}}
\toprule
\textbf{Method} & \textbf{Formula} & \textbf{Use Case} \\
\midrule
Hoeffding & $n \geq \frac{\ln(2/\delta)}{2\epsilon^2}$ & General, conservative \\
Chernoff & $n \geq \frac{3\ln(2/\delta)}{\epsilon^2 p}$ & Small probabilities \\
Clopper-Pearson & Exact binomial & High confidence \\
\bottomrule
\end{tabular}
\end{table}

\begin{lstlisting}[language=Python,caption=Shot Optimization]
from qpu_estimator.estimator import EstimationConfig

config = EstimationConfig(
    shot_method="hoeffding",  # or "chernoff", "clopper-pearson"
    target_precision=0.01,
    confidence=0.95,
)
\end{lstlisting}

%=============================================================================
\section{QML Estimation}
%=============================================================================

\subsection{VQC Training}

Estimates total resources for variational quantum classifier training:

\begin{lstlisting}[language=Python,caption=VQC Training Estimation]
from qpu_estimator import QMLEstimator
from qpu_estimator.qml_estimator import VQCEstimationConfig

qml = QMLEstimator()
config = VQCEstimationConfig(
    num_parameters=4,
    num_training_epochs=100,
    gradient_method="spsa",  # or "parameter-shift", "finite-diff"
)

report = qml.estimate_vqc_training(
    circuit, "ibm_kingston", config
)
print(f"Total shots: {report.optimal_shots:,}")
print(f"Final fidelity: {report.estimated_fidelity:.4f}")
\end{lstlisting}

\subsection{MAML/QMAML}

Estimates meta-learning resources:

\begin{lstlisting}[language=Python,caption=MAML Estimation]
from qpu_estimator.qml_estimator import MAMLEstimationConfig

config = MAMLEstimationConfig(
    num_tasks=10,
    inner_loop_steps=5,
    outer_loop_steps=20,
)
report = qml.estimate_maml(circuit, "ibm_kingston", config)
\end{lstlisting}

\subsection{Transfer Learning}

Estimates quantum transfer learning overhead:

\begin{lstlisting}[language=Python,caption=Transfer Learning Estimation]
from qpu_estimator.qml_estimator import TransferLearningConfig

config = TransferLearningConfig(
    classical_layers=3,
    quantum_layers=2,
    frozen_classical=True,
    fine_tune_epochs=20,
)
report = qml.estimate_transfer_learning(
    quantum_circuit, "ibm_kingston", config
)
\end{lstlisting}

%=============================================================================
\section{Training Time Estimation}
%=============================================================================

\subsection{Full Experiment Time}

Estimates total wall-clock time including:

\begin{itemize}[leftmargin=2em]
    \item Forward/backward passes per batch
    \item Quantum evaluation time
    \item Classical overhead (data loading, optimizer)
    \item API communication overhead
    \item Queue time (heuristic by plan priority)
\end{itemize}

\begin{lstlisting}[language=Python,caption=Training Time Estimation]
from qpu_estimator import TrainingTimeEstimator

time_est = TrainingTimeEstimator()
result = time_est.estimate_full_experiment(
    circuit=circuit,
    backend_name="ibm_kingston",
    n_samples=3772,
    batch_size=32,
    n_epochs=100,
    gradient_method="spsa",
    n_params=4,
    plan_priority="open",  # or "paygo", "premium"
)

print(f"Total time: {result['total_time_hours']:.2f} hours")
print(f"Queue time: {result['queue_time_ms']/3600000:.2f} hours")
\end{lstlisting}

\subsection{Session Mode}

IBM Runtime Session mode queues once and runs all circuits in a single session:

\begin{lstlisting}[language=Python,caption=Session Mode Estimation]
session = time_est.estimate_with_session(
    circuit, "ibm_kingston",
    n_samples=3772, batch_size=32, n_epochs=100,
    plan_priority="open",
)
print(f"Speedup: {session['speedup_vs_standard']:.1f}x")
\end{lstlisting}

\subsection{Plan Priority Comparison}

\begin{table}[htbp]
\centering
\caption{Queue Time Heuristics by Plan}
\begin{tabular}{@{}lll@{}}
\toprule
\textbf{Plan} & \textbf{Queue Time} & \textbf{Use Case} \\
\midrule
open & 30 min/job & Free, limited priority \\
paygo & 5 min/job & Pay-as-you-go \\
premium & 1 min/job & Enterprise priority \\
\bottomrule
\end{tabular}
\end{table}

%=============================================================================
\section{Export Utilities}
%=============================================================================

\subsection{LaTeX Export}

\begin{lstlisting}[language=Python,caption=LaTeX Table Export]
from qpu_estimator import ReportExporter

reports = estimator.compare_backends(circuit, backends)
latex = ReportExporter.to_latex_table(
    reports, caption="Backend Comparison"
)
print(latex)
\end{lstlisting}

Output:
\begin{lstlisting}[language=TeX,caption=Generated LaTeX]
\begin{table}[htbp]
\centering
\caption{Backend Comparison}
\begin{tabular}{lcccccc}
\toprule
Backend & Depth & Time (ms) & Shots & Fidelity & Credits & SWAPs \\
\midrule
ibm\_kingston & 7 & 0.003 & 18445 & 0.8831 & 0.0001 & 0 \\
\bottomrule
\end{tabular}
\end{table}
\end{lstlisting}

\subsection{CSV and Markdown}

\begin{lstlisting}[language=Python,caption=CSV and Markdown Export]
csv = ReportExporter.to_csv(reports)
md = ReportExporter.to_markdown(reports)
\end{lstlisting}

%=============================================================================
\section{Plugin System}
%=============================================================================

\subsection{Creating a Custom Backend Plugin}

\begin{lstlisting}[language=Python,caption=Custom Backend Plugin]
from qpu_estimator.plugin import BackendPlugin, PluginRegistry
from qpu_estimator.models import BackendProfile

class IonQPlugin(BackendPlugin):
    @property
    def name(self) -> str:
        return "ionq_harmony"
    
    def get_profile(self) -> BackendProfile:
        return BackendProfile(
            name=self.name,
            num_qubits=11,
            basis_gates=["gpi", "gpi2", "ms"],
            coupling_map=[[i, j] for i in range(11) for j in range(i+1, 11)],
            t1_times_us=[10000.0] * 11,
            t2_times_us=[10000.0] * 11,
            single_qubit_error=[0.0001] * 11,
            two_qubit_error=[0.003] * 55,
            readout_error=[0.005] * 11,
            gate_times_ns={"gpi": 100, "gpi2": 100, "ms": 500},
        )

# Register
registry = PluginRegistry()
registry.register(IonQPlugin())
\end{lstlisting}

%=============================================================================
\section{Command-Line Interface}
%=============================================================================

\subsection{Basic Usage}

\begin{lstlisting}[language=bash,caption=CLI Examples]
# Single backend estimation
qpu-estimate circuit.qpy --backend ibm_kingston

# Multi-backend comparison
qpu-estimate circuit.qpy --compare ibm_fez ibm_marrakesh ibm_kingston

# JSON output
qpu-estimate circuit.qpy --backend ibm_kingston --json
\end{lstlisting}

\subsection{Output Format}

\begin{lstlisting}[language=bash,caption=Human-Readable Output]
==================================================
Backend: ibm_kingston
==================================================
Circuit qubits:        2
Total gates:           5
Single-qubit gates:    1
Two-qubit gates:       1
Original depth:        3
Transpiled depth:      7
SWAP count:            0
Execution time:        0.003 ms
Optimal shots:         18445
Estimated fidelity:    0.8831
Estimated credits:     0.0001

Notes:
  - Noise-aware fidelity model
  - Shot optimization method: hoeffding
==================================================
\end{lstlisting}

%=============================================================================
\section{Validation Results}
%=============================================================================

\subsection{IBM Quantum Live Execution}

Job \texttt{d8vo9femvj5c73ei7adg} executed on \texttt{ibm\_kingston}:

\begin{table}[htbp]
\centering
\caption{Framework Estimation vs. Real Execution}
\begin{tabular}{@{}lcc@{}}
\toprule
\textbf{Metric} & \textbf{Estimated} & \textbf{Real} \\
\midrule
Backend & ibm\_kingston & ibm\_kingston \\
Depth (transpiled) & 7 & 7 \\
SWAPs & 0 & 0 \\
Shots & 18445 & 100 \\
Fidelity & 0.8831 & 0.9900 \\
Execution time & 0.003 ms & $\sim$0.003 ms \\
Queue time & Not estimated & $\sim$5 min \\
\bottomrule
\end{tabular}
\end{table}

\subsection{Error Analysis}

\begin{itemize}[leftmargin=2em]
    \item \textbf{Depth:} Exact (0\% error)
    \item \textbf{Fidelity:} Conservative underestimate ($-0.11$). The framework is designed to be pessimistic for safe planning.
    \item \textbf{Queue time:} Not yet implemented (backlog item)
\end{itemize}

\subsection{Simulator Validation (Aer)}

Using \texttt{NoiseModel.from\_backend(ibm\_kingston)}:

\begin{table}[htbp]
\centering
\caption{Framework vs. Aer Simulator}
\begin{tabular}{@{}lccc@{}}
\toprule
\textbf{Circuit} & \textbf{Estimated} & \textbf{Simulated} & \textbf{Error} \\
\midrule
Bell state & 0.8831 & 0.9904 & 0.1073 \\
GHZ-3 & 0.8558 & 0.9820 & 0.1262 \\
GHZ-4 & 0.8069 & 0.9746 & 0.1677 \\
GHZ-5 & 0.7325 & 0.9616 & 0.2291 \\
\bottomrule
\end{tabular}
\end{table}

%=============================================================================
\section{Testing}
%=============================================================================

\subsection{Test Suite}

\begin{lstlisting}[language=bash,caption=Running Tests]
pytest tests/ -v
\end{lstlisting}

\begin{table}[htbp]
\centering
\caption{Test Coverage by Phase}
\begin{tabular}{@{}lll@{}}
\toprule
\textbf{Phase} & \textbf{Tests} & \textbf{Status} \\
\midrule
MVP (analyzer, profiler, transpiler) & 8 & Passing \\
IBM Integration (live profiling) & 9 & Passing \\
Advanced Estimation (noise, shots) & 5 & Passing \\
QML Specialization (VQC, MAML) & 5 & Passing \\
Ecosystem (export, plugins) & 6 & Passing \\
Training Time Estimation & 4 & Passing \\
\midrule
\textbf{Total} & \textbf{37/39} & \textbf{2 skipped (need env token)} \\
\bottomrule
\end{tabular}
\end{table}

%=============================================================================
\section{Future Work}
%=============================================================================

\begin{itemize}[leftmargin=2em]
    \item \textbf{Queue time estimation:} Historical data + heuristics
    \item \textbf{IBM credit pricing table:} Manual + scrape fallback
    \item \textbf{CI integration:} Fail if estimated cost exceeds budget
    \item \textbf{Caching:} Backend profiles to avoid repeated API calls
    \item \textbf{Parallel estimation:} Across multiple backends simultaneously
    \item \textbf{Historical tracking:} Store estimations vs. actual execution metrics
\end{itemize}

%=============================================================================
\section{License and Contributing}
%=============================================================================

QPU-Estimator is released under the MIT License.

\begin{itemize}[leftmargin=2em]
    \item Repository: \url{https://github.com/danimap27/qpu-estimator}
    \item Issues: \url{https://github.com/danimap27/qpu-estimator/issues}
    \item Author: Daniel Mart\'{i}n P\'{e}rez (danimap27@gmail.com)
\end{itemize}

\end{document}
